\subsection{Four objects with cyclic support at nineteen arrows}

\begin{lemma}[Proved: a singleton long hom-set across a one-arrow minimal ideal]
\label{CP2:A-singleton-long-minimal-pair-retract}
Let a four-object category have condensation $u<S<v$, where $S$
is a strong pair whose minimal ideal has one arrow in each ordered
hom-set. If $|C(u,v)|=1$, a projective-plane category of this form
admits an actual arrow-reducing projective-plane retract or homotopy
equivalence.
\end{lemma}
\begin{proof}
Let $e_i$ be the unique minimal endomorphism at each object of $S$.
For every internal arrow $f:i\to j$,
$e_jf=fe_i$, since both sides are the unique arrow of $K(i,j)$.
On the connecting profunctor for the cut $\{u\}\to S\cup\{v\}$,
send an entrance $p:u\to i$ to $e_ip$ and fix the unique long
arrow. This is natural for all arrows of $S$ by the displayed
identity, and for the exits to $v$ because their composites lie
in the singleton long hom-set. It also commutes with both outside
corner actions. Thus it is an idempotent natural transformation
of the actual connecting profunctor.

If its image is proper, the cut-retract theorem gives an actual
projective-plane retract with fewer arrows. Otherwise all external
entrances are fixed by the minimal idempotents. There is one minimal
row per object and trivial local group, so the fixed-entrance
component collapse of Lemma~\ref{CP2:A-high-one-sided} is an actual
homotopy equivalence replacing $S$ by a point, strictly reducing
arrows. These are the two asserted alternatives.
\end{proof}

\begin{lemma}[Proved: annihilated augmentations at a strong pair of size at most eight]
\label{CP2:A-eight-pair-annihilated-augmentation-Tor}
Let a category have condensation $u<S<v$, where $|S|\le8$,
neither object of $S$ is a retract, and both outside corner
algebras are semisimple over $\mathbb Q$. Let $V,W$ be the
entrance and exit augmentation modules. If the minimal ideal
$K$ of $S$ annihilates both modules, bounded rational homology
of the whole category implies $H_4=0$.
\end{lemma}
\begin{proof}
The minimal group is trivial, since a nontrivial one would cost
at least $2\cdot2\cdot2+2=10$ internal arrows. Write
$R=\mathbb QS$ and $A=R/\mathbb QK$. The ideal is projective
on both sides and has the form $Re\otimes eR$. Since $We=0$,
change of rings identifies $\operatorname{Tor}^R(W,V)$ with
$\operatorname{Tor}^A(W,V)$, as in
Lemma~\ref{CP2:A-small-pair-augmentation-Tor}.

Apply exact outside orbit functors to make both extreme corners
trivial. The new augmentation modules are exact coinvariant
quotients of the original ones, so are still annihilated by $K$.
These actual functors preserve rational homology and the cohomology
ring. The middle-bar formula identifies the new fourth homology
with the second Tor over $A$, and the higher groups with the same
Tor groups shifted by two. The quotient $A$ has dimension at most
four and has two nonzero orthogonal object identities, so is
nonlocal. Lemma~\ref{A:nonlocal-four} implies that a nonzero
second Tor would give unbounded higher Tor. This contradicts the
transferred boundedness and proves the assertion.
\end{proof}

\begin{theorem}[Proved: cyclic four-object support through nineteen]
\label{CP2:A-strong-four-nineteen}
A finite category on four objects with cyclic support and at most
nineteen arrows cannot have classifying space of type $\CP^2$.
\end{theorem}
\begin{proof}
Any actual arrow-reducing projective-plane retract or homotopy
equivalence contradicts the complete eighteen-arrow exclusion.
Assume in particular that there is no object retract. A wholly
strong four-object category then costs at least twenty arrows.
A proper three-object strong component costs at least twelve;
its interface costs at least three, so the remaining corner has
order at most four. The component has finite-group type, since
otherwise its internal minimum is nineteen. Both blocks have
zero second rational cohomology. Two strong pairs likewise have
an interface of at least four arrows; neither can have nonzero
second cohomology, since that would require at least
$11+6+4=21$ arrows. The two-block lemma excludes both cases.

There remains one strong pair $S$, of size $q\ge6$, with a
two-object complement $D$.

\emph{An extreme pair.} By opposition suppose $S$ is a source.
Write $h$ for the total interface cardinality. Connectedness gives
$q\le14$ and $h\ge2$. If both block second cohomology groups
vanish, use the two-block lemma.

Suppose first that $H^2(BD;\mathbb Q)\ne0$, so $|D|\ge6$.
At $q=6,7$ there is one minimal row per object and the source
collapse applies. At $q=8$ the same applies unless the total
rows and columns are $(3,2)$, in which case the pair is pure
and its one-row object admits incoming one-sided deletion.
We may therefore assume $9\le q\le11$.

If $D$ is disconnected, its positive-second-homology minimum is
six and $h\ge4$. The budget forces $q=9$, $|D|=6$, $h=4$.
Every interface stalk is singleton, so the whole category is the
categorical join of $S$ and $D$ and its positive products vanish.
Suppose $D$ is connected. If $S$ meets both objects, the same
budget forces $q=9$, $|D|=6$, $h=4$, again a full singleton join.
Otherwise it meets one object $w$ only; the other object points
to $w$, because the reverse orientation would force another
interface. The nerves are the union of $E=S\cup\{w\}$ and
$D$, meeting in the corner $N$ at $w$.

Every connected two-object category with acyclic support and size
at most eight has zero rational second products. Indeed, if both corners have
order at most four, their second groups vanish and the two-block
lemma applies. Otherwise either the stalk is singleton or it has
two arrows; the join or unrestricted two-cross theorem applies.
Moreover a positive second group at size at most eight forces
both corners to have order at most five, and at size at most
seven forces both to have order at most four.

If $h=2$, then $|D|\le8$, and $E$ is a singleton join, whose
positive products vanish. Its intersection corner $N$ has order
at most five and hence zero third rational homology. Fourth
restriction to $E,D$ is injective, so all global second products
vanish. If $h\ge3$, then $|D|\le7$ and $N$ is rationally
acyclic. Mayer--Vietoris identifies
$H^2(BC)=H^2(BE)\oplus H^2(BD)$. The second summand is nonzero
and the total has dimension one, so $H^2(BE)=0$. The global
products restrict to zero on $E$ and, by the preceding two-object
obstruction, on $D$. Fourth restriction is injective because
$H^3(BN)=0$. This completes the positive-complement branch.

Now assume $H^2(BD;\mathbb Q)=0$ and
$H^2(BS;\mathbb Q)\ne0$. Then $11\le q\le14$.
If $D$ is disconnected, $h\ge4$ and $|D|\le4$.
The case $|D|=4$ forces $q=11,h=4$, a full singleton join.
Otherwise both leaf corners have order at most two. Their exact
orbit functors and the usual leaf-merging comparison give a
bounded ordinary act cone over $S$, preserving second products.
Here $q\le13$, so the total minimal row and column counts are
both three. Every local row count is at most two, and the
binary projective-cover theorem kills those products.

For connected $D$, if $S$ meets both outside objects, the budget
gives $h\le5$. At least one of the two target objects therefore
receives singleton stalks from both objects of $S$. The inclusion
of $S$ is nullhomotopic through that target, and the other block
has zero second cohomology; apply the two-block lemma.
Otherwise $S$ meets only the final object $w$ of the outside edge.
Again write $E=S\cup\{w\}$ and $N=\operatorname{End}(w)$.
The order of $N$ is at most $19-q-h-2\le4$. The actual nerve
union has rationally acyclic intersection $BN$, so $E$ is bounded,
$\dim H^2(BE)=1$, and fourth restriction to $E,D$ is injective.
The case $h=2$ is a singleton join. For $h\ge3$, $q\le13$
and each local minimal row count of $S$ is at most two. If
$|N|\le2$, exact sink orbits and the binary projective-cover
theorem apply. If $|N|\ge3$, the budget forces $h=3$, so both
stalks are binary and the small-binary sink lemma applies.
The restriction to $D$ is already zero in degree two. This
exhausts the extreme pair.

\emph{A nonextreme pair.} The condensation is $u<S<v$.
Let $b,c$ be the endpoint corner orders, let $z=|C(u,v)|$, and
let $\alpha_i,\beta_i$ be the entrance and exit sizes at the
two objects of $S$. Put
\[
 P=\sum_i(\alpha_i-1)+\sum_i(\beta_i-1),\qquad
 \delta=(b-1)+(c-1)+(z-1).
\]
Then
\[
 q+7+P+\delta\le19,\qquad q\le12,
 \qquad P+\delta\le12-q.
\]
Whenever all four entrance and exit stalks are singleton, their
composites give a single long arrow $z_0$, independent of the
chosen object of $S$ by strong connectedness. It is fixed by both
endpoint actions. If $z>1$, fixing the entrance stalks and sending
every long arrow to $z_0$ gives a proper natural cut retraction;
if $z=1$, the two outside corners and $S$ form a categorical join.
Thus $P=0$ is excluded in every case.

For $q\ge9$ and $z\le2$, both endpoint corners have order at
most four. Failure of the binary-stalk test at both ends requires
an entrance and an exit of size at least three, costing $P\ge4$,
whereas $P\le3$. One small-binary endpoint therefore applies.
For $q\ge9$ and $z\ge3$, the budget gives $P\le1$ and, after
the already excluded $P=0$ case, both endpoint corners have order
at most two. Exact outside orbits make them trivial. One entire
entrance or exit side is singleton, so its augmentation module
is zero and the middle-bar formula gives $H_4=0$.

At $q=8$ and $z\le2$, if $P\le3$ then the same binary-endpoint
argument applies: $P>0$ makes both endpoint corners at most four.
If $P=4$, the budget forces $b=c=z=1$. Avoiding a binary end
then forces entrance sizes $(3,1)$ and exit sizes $(3,1)$, in
either order. Both augmentation modules are supported on one
object. Every minimal endomorphism factors through the other
object, so $K$ annihilates both. Apply the preceding Tor lemma.
For $q=8,z\ge3$ and $P>0$, the endpoint corners have order at
most two and $P\le2$. If one entire side is singleton, its
augmentation is zero. Otherwise $P=2$ and each augmentation is
supported on one object, so the same Tor lemma applies.

It remains $q=6$ or $7$. Their minimal ideal has one arrow in
each ordered hom-set, and its local idempotents are corner zeros.
For $z=1$, the singleton-long retract lemma excludes the category.
Assume $z\ge2$. If any of the four local incoming/outgoing
fixed-arrow conditions holds, one-sided object deletion applies;
the corner zero supplies its unit hypothesis. Assume all four fail.

At $q=6$, all internal cross arrows are minimal, so each of the
four external stalks must be nonsingleton and $P\ge4$.
For $z=2$ one has $P\le5$ and both endpoint corners are at most
two. Failure of binary tests at both ends would require $P\ge6$,
so a binary endpoint applies. For $z\ge3$, the budget forces
$z=3$, $P=4$, $b=c=1$ and all four stalks of size two. Each
minimal idempotent is a nonidentity idempotent on each relevant
binary stalk, hence is constant and kills its augmentation.
Every minimal arrow factors through these idempotents, so $K$
annihilates both augmentation modules. The Tor lemma applies.

At $q=7$, the single extra internal arrow can affect at most two
fixed-arrow conditions, hence $P\ge2$. If $z=2$, then $P\le4$
and both endpoint corners are at most three. A binary endpoint
applies unless $P=4$, which forces $b=c=1$ and one ternary and
one singleton stalk on each side. Their supported augmentations
are killed by $K$, so the Tor lemma applies.

Finally let $q=7,z\ge3$. Then $2\le P\le3$ and both endpoint
corners have order at most two. If an entire side is singleton,
the middle-bar group is zero. Otherwise each augmentation has
dimension at most two. If it is supported at one object, $K$
kills it by factorization through the other. The only remaining
possibility on either side is two binary stalks. At least one
object has all internal incoming cross arrows fixed by its
minimal idempotent: the single extra arrow can have at most one
distinct target. Since all four full conditions fail, that
idempotent is nonidentity on its binary entrance, hence acts
as zero on its augmentation. The unique minimal arrow from the
other object into this one is then zero on augmentations, and
factorization through it shows that the other minimal idempotent
is zero too. Thus all of $K$ kills the entrance augmentation.
The outgoing argument is dual and gives the same conclusion for
the exit augmentation. Exact outside coinvariants preserve this
annihilation, and the Tor lemma proves $H_4=0$. Every nonextreme
budget has now been excluded.
\end{proof}
